% BIE = boundary integral equation.
\section{Introduction}\label{sec:intro}
Moving an artificial vertical boundary inside a homogeneous strip of a
photonic crystal leaves all rods in place. It can nevertheless move a pole
of the half-strip Dirichlet-to-Neumann (DtN) map directly onto a guided
wavenumber. A solver that locates zeros of a DtN determinant then faces two
distinct problems. An isolated pole can give a small ratio of extreme
singular values and pass a numerical acceptance test. If the pole and a
mode coincide, division by the singular Dirichlet trace can cancel the
zero itself. In the latter case even an exact evaluation of the reduced
determinant need not reveal the mode.

We consider a line defect between periodic dielectric half-strips. They
share the period along the defect and may differ in the transverse period,
rod shape and material. At fixed tangential Bloch wavenumber, an eigenvalue
in a common gap of the surrounding crystals has a field decaying into
both half-strips \cite{Fliss2013,KlindworthSchmidtFliss2014}. Our aim is to
compute such wavenumbers and fields by discretizing the material interfaces.
The method uses a boundary integral equation (BIE) for the response of a
cell, Rayleigh amplitudes to exchange fields between adjacent cells, and
Floquet modes to represent decay at infinity.

\paragraph{Transparent boundaries and existing solvers}
Cell-based transparent conditions reduce the unbounded problem to the
defect region \cite{JolyLiFliss2006,Fliss2013}.
Klindworth, Schmidt and Fliss \cite{KlindworthSchmidtFliss2014} treated
different periodic claddings by finite elements. Fliss, Klindworth and
Schmidt \cite{FlissKlindworthSchmidt2015} distinguished global half-strip
Dirichlet exceptional values from local cell exceptional values. Their
Robin-to-Robin (RtR) formulation handles these values through Robin data.
It provides a natural comparison for a method that already has access to
both components of the trace. In the present formulation the full
half-strip data are available directly from the decaying Floquet subspace.
We retain its coefficients and impose membership in that subspace on each
side of the defect.

\paragraph{Scattering matrices and boundary integral determinants}
Gralak, Enoch and Tayeb \cite{GralakEnochTayeb2002} compute Bloch modes
from a layer scattering matrix. Botten et al.\ \cite[eq.~(38)]{BottenEtAl2004}
write continuity at an interface between two photonic crystals as equality
of the complete plane-wave states expanded in the two Bloch bases. They
subsequently derived a determinant equation for defect modes in cylinder
arrays \cite{BottenEtAl2006}. Their W3 guide supplies an external benchmark
below. In a BIE for forced scattering by periodic layers, Cho and Barnett
\cite[eqs.~(47) and~(53)]{ChoBarnett2015} retain the outgoing Rayleigh
amplitudes together with the layer densities and match the field and its
normal derivative. Here the exterior states are the decaying Floquet
subspaces of periodic half-strips.
Barnett and Greengard \cite{BarnettGreengard2010} showed how a change of
integral representation avoids empty-cell resonances in a doubly periodic
band-structure calculation. Zhao and Barnett \cite{ZhaoBarnett2015} used an
analytic Fredholm determinant and polynomial interpolation to compute
Dirichlet eigenvalues of a bounded domain. Here the representation
singularity arises when the decaying half-strip data are written as a graph
over their Dirichlet traces. For a direct subspace determinant, orthonormal bases returned by independent
generalized Schur decompositions need not vary continuously with frequency.

\paragraph{The discrete guided-mode problem}
We apply direct modal matching to the defect-cell BIE and the two decaying
Floquet subspaces. We prove the dimension count for the flux-conserving
truncated cell and identify the kernel of the resulting bordered matrix.
For this discretization, block elimination gives the precise relation
between the full matching determinant and the DtN determinant. The factors include the center density block, both Dirichlet
trace blocks and the Rayleigh-to-Cauchy coordinate changes. Their orders
explain when Dirichlet elimination can obscure a mode in this solver.
A locally normalized basis then yields an analytic
determinant wherever the cell assembly and the subspace separation are
regular. The same trace data also give the discrete RtR formulation.

To test the mechanism we compute the same waveguide with the artificial
vertical boundary in three positions. One position puts a Dirichlet pole
near the mode and another aligns them to numerical precision. We compare
the returned frequencies, the status of the whole search window and jumps
of the reconstructed fields. An empty defect between unequal-period
half-strips, the W3 guide and a finite element refinement provide separate
checks. The mathematical kernel statements concern the discrete matching
equations. The experiments assess how their roots approximate physical
wavenumbers.

\Cref{sec:model} defines the waveguide and the frequency window.
\Cref{sec:discretization,sec:leads} construct the cell data and the decaying
spaces. \Cref{sec:closure} derives the bordered determinant and its relation
to the DtN and RtR descriptions. \Cref{sec:algorithm,sec:experiments}
describe the search, field reconstruction and numerical tests.
